\chapter{3D Cartesian Coordinates and Direction Cosines (Pages 15--24)}
\label{ch:coordinates_dc_dr}

% =========================================================================
% PAGE 15
% =========================================================================
\section{Page 15: Introduction to 3D Cartesian Coordinates and Section Formulae}
\label{sec:page15}

\begin{conceptbox}[Manuscript Overview: Page 15]
Page 15 transitions the course from two-dimensional polar geometry to three-dimensional analytical geometry. It introduces the rectangular Cartesian coordinate frame in $\mathbb{R}^3$, the definition of points $P(x_1, y_1, z_1)$ and $Q(x_2, y_2, z_2)$, the section formula dividing a segment in the ratio $m_1 : m_2$, the midpoint formula, and the concept of changing coordinates via translation of axes.
\end{conceptbox}

\subsection{The Spatial Cartesian Coordinate System}
To represent spatial configurations analytically, three mutually perpendicular straight lines are chosen that intersect at a single point $O$, known as the \textbf{origin}:
\begin{itemize}[noitemsep]
    \item $X'OX$: The \textbf{$x$-axis} (abscissa axis).
    \item $Y'OY$: The \textbf{$y$-axis}.
    \item $Z'OZ$: The \textbf{$z$-axis} (applicate axis).
\end{itemize}
These three lines define three mutually perpendicular planes, known as the \textbf{coordinate planes}:
\begin{enumerate}[noitemsep]
    \item The \textbf{$xy$-plane}: $z = 0$, spanned by the $x$- and $y$-axes.
    \item The \textbf{$yz$-plane}: $x = 0$, spanned by the $y$- and $z$-axes.
    \item The \textbf{$zx$-plane}: $y = 0$, spanned by the $z$- and $x$-axes.
\end{enumerate}
These planes divide the whole of three-dimensional Euclidean space into $8$ compartments, termed \textbf{octants}. The signs of the triad of real numbers $(x, y, z)$ determine the octant in which any point $P$ resides.

\begin{center}
\begin{tikzpicture}[scale=1.1, >=Stealth, 3d view={120}{25}]
    % Axes
    \draw[->, thick] (0,0,0) -- (4.8,0,0) node[anchor=north east]{$x$};
    \draw[->, thick] (0,0,0) -- (0,5.2,0) node[anchor=north west]{$y$};
    \draw[->, thick] (0,0,0) -- (0,0,4.2) node[anchor=south]{$z$};

    % Origin
    \coordinate (O) at (0,0,0);
    \node[below left] at (O) {$O(0,0,0)$};

    % Coordinates of P
    \def\px{3.2}
    \def\py{4.0}
    \def\pz{2.8}

    \coordinate (P) at (\px,\py,\pz);
    \coordinate (A) at (\px,0,0);
    \coordinate (B) at (0,\py,0);
    \coordinate (C) at (0,0,\pz);
    \coordinate (M) at (\px,\py,0); % xy projection
    \coordinate (N) at (0,\py,\pz); % yz projection
    \coordinate (L) at (\px,0,\pz); % zx projection

    % Box edges
    \draw[dashed, gray!80] (A) -- (M) -- (B);
    \draw[dashed, gray!80] (O) -- (M);
    \draw[dashed, secondary, thick] (M) -- (P) node[midway, right] {$z$};
    \draw[dashed, gray!80] (C) -- (L) -- (A);
    \draw[dashed, gray!80] (C) -- (N) -- (B);
    \draw[dashed, gray!80] (L) -- (P);
    \draw[dashed, gray!80] (N) -- (P);

    % Distance OP
    \draw[very thick, primary] (O) -- (P) node[midway, above left] {$r = OP$};

    % Points
    \fill (O) circle (2pt);
    \fill[accent] (P) circle (2.5pt) node[above right] {$P(x,y,z)$};
    \fill[secondary] (M) circle (1.8pt) node[below right] {$M(x,y,0)$};
\end{tikzpicture}
\end{center}

\subsection{The 3D Distance Formula}
Let $P(x_1, y_1, z_1)$ and $Q(x_2, y_2, z_2)$ be two arbitrary points in space. Construct a rectangular cuboid whose opposite vertices are $P$ and $Q$, with faces parallel to the coordinate planes. The lengths of the edges of this cuboid are:
\begin{equation}
    \Delta x = |x_2 - x_1|, \quad \Delta y = |y_2 - y_1|, \quad \Delta z = |z_2 - z_1|.
\end{equation}
By double application of the Pythagorean theorem:
\begin{equation}
    PQ^2 = (\Delta x)^2 + (\Delta y)^2 + (\Delta z)^2 = (x_2 - x_1)^2 + (y_2 - y_1)^2 + (z_2 - z_1)^2.
\end{equation}
Thus, the Euclidean distance between $P$ and $Q$ is:
\begin{equation}
    \boxed{PQ = \sqrt{(x_2 - x_1)^2 + (y_2 - y_1)^2 + (z_2 - z_1)^2}}
\end{equation}
In particular, the distance of any point $P(x, y, z)$ from the origin $O(0,0,0)$ is:
\begin{equation}
    OP = \sqrt{x^2 + y^2 + z^2}.
\end{equation}

\subsection{Section Formula (Internal and External Division)}
\begin{theorembox}[Section Formula in Three Dimensions]
Let $P(x_1, y_1, z_1)$ and $Q(x_2, y_2, z_2)$ be two given points. Let $R(x, y, z)$ be a point dividing the directed line segment $PQ$ in the given ratio $m_1 : m_2$.
\begin{enumerate}
    \item \textbf{Internal Division}: If $R$ lies between $P$ and $Q$,
    \begin{equation}
        \boxed{x = \frac{m_1 x_2 + m_2 x_1}{m_1 + m_2}, \quad y = \frac{m_1 y_2 + m_2 y_1}{m_1 + m_2}, \quad z = \frac{m_1 z_2 + m_2 z_1}{m_1 + m_2}}
    \end{equation}
    \item \textbf{External Division}: If $R$ lies on the line produced outside the segment $PQ$,
    \begin{equation}
        \boxed{x = \frac{m_1 x_2 - m_2 x_1}{m_1 - m_2}, \quad y = \frac{m_1 y_2 - m_2 y_1}{m_1 - m_2}, \quad z = \frac{m_1 z_2 - m_2 z_1}{m_1 - m_2}} \quad (m_1 \neq m_2).
    \end{equation}
\end{enumerate}
\end{theorembox}

\begin{proof}
Draw perpendiculars $PP', QQ', RR'$ from $P, Q, R$ respectively to the $xy$-plane. The feet of these perpendiculars $P', Q', R'$ lie in the $xy$-plane and have coordinates $(x_1, y_1, 0)$, $(x_2, y_2, 0)$, and $(x, y, 0)$.
Since $PP', RR', QQ'$ are all perpendicular to the $xy$-plane, they are parallel to one another and lie in a common plane.
Through $R$, draw a line parallel to $P'Q'$ meeting $PP'$ at $S$ and $QQ'$ at $T$.
From the elementary geometry of similar triangles $\triangle PSR \sim \triangle RTQ$:
\begin{equation}
    \frac{PR}{RQ} = \frac{m_1}{m_2} = \frac{z - z_1}{z_2 - z}.
\end{equation}
Cross-multiplying gives:
\begin{equation}
    m_1(z_2 - z) = m_2(z - z_1) \implies (m_1 + m_2)z = m_1 z_2 + m_2 z_1 \implies z = \frac{m_1 z_2 + m_2 z_1}{m_1 + m_2}.
\end{equation}
By projecting onto the $yz$-plane and $zx$-plane in an identical manner, we obtain:
\begin{equation}
    x = \frac{m_1 x_2 + m_2 x_1}{m_1 + m_2}, \quad y = \frac{m_1 y_2 + m_2 y_1}{m_1 + m_2}.
\end{equation}
For external division, replacing $m_2$ by $-m_2$ yields the required result.
\end{proof}

\begin{notebox}[Important Corollaries of the Section Formula]
\begin{itemize}[noitemsep]
    \item \textbf{Midpoint Formula}: When $m_1 : m_2 = 1 : 1$, the midpoint $M$ of segment $PQ$ is:
    \begin{equation}
        M = \left(\frac{x_1 + x_2}{2}, \; \frac{y_1 + y_2}{2}, \; \frac{z_1 + z_2}{2}\right).
    \end{equation}
    \item \textbf{Centroid of a Triangle}: If vertices are $A(x_1, y_1, z_1), B(x_2, y_2, z_2), C(x_3, y_3, z_3)$, the medians intersect at centroid $G$:
    \begin{equation}
        G = \left(\frac{x_1+x_2+x_3}{3}, \; \frac{y_1+y_2+y_3}{3}, \; \frac{z_1+z_2+z_3}{3}\right).
    \end{equation}
    \item \textbf{Centroid of a Tetrahedron}: If vertices are $A_i(x_i, y_i, z_i)$ for $i=1,2,3,4$, the centroid is:
    \begin{equation}
        G = \left(\frac{\sum x_i}{4}, \; \frac{\sum y_i}{4}, \; \frac{\sum z_i}{4}\right).
    \end{equation}
\end{itemize}
\end{notebox}

% =========================================================================
% PAGE 16
% =========================================================================
\section{Page 16: Transformation of Coordinates --- Translation of Axes}
\label{sec:page16}

\begin{conceptbox}[Manuscript Overview: Page 16]
Page 16 explores the kinematics of coordinate frames in space. It focuses on the translation of axes: shifting the origin from $O(0,0,0)$ to a new point $O'(h, k, l)$ while maintaining the parallel orientation of all three coordinate axes.
\end{conceptbox}

\subsection{Transformation Equations for Translation}
Let $OX, OY, OZ$ be the original coordinate axes with origin $O(0,0,0)$.
Let the origin be shifted to a new point $O'$ whose coordinates with respect to the original system are $(h, k, l)$.
Let through $O'$ three new axes $O'X', O'Y', O'Z'$ be drawn parallel respectively to $OX, OY, OZ$ and directed in the same sense.

\begin{center}
\begin{tikzpicture}[scale=1.1, >=Stealth]
    % Old axes
    \draw[->, thick, gray] (0,0) -- (5.5,0) node[right] {$X$};
    \draw[->, thick, gray] (0,0) -- (0,4.5) node[above] {$Y$};
    \node[below left] at (0,0) {$O(0,0,0)$};

    % New origin O'
    \coordinate (Oprime) at (2.2, 1.4);
    \fill[primary] (Oprime) circle (2pt) node[below right] {$O'(h,k,l)$};
    \draw[dashed, gray] (2.2,0) node[below] {$h$} -- (Oprime) -- (0,1.4) node[left] {$k$};

    % New axes
    \draw[->, very thick, primary] (Oprime) -- ++(3.8,0) node[right] {$X'$};
    \draw[->, very thick, primary] (Oprime) -- ++(0,3.2) node[above] {$Y'$};

    % Point P
    \coordinate (P) at (4.5, 3.6);
    \fill[accent] (P) circle (2.2pt) node[above right] {$P(x,y,z) \equiv P(x',y',z')$};

    % Vector arrows
    \draw[->, thick, secondary] (0,0) -- (Oprime) node[midway, below right] {$\vec{r}_0$};
    \draw[->, thick, accent] (Oprime) -- (P) node[midway, below right] {$\vec{r}'$};
    \draw[->, very thick, primary] (0,0) -- (P) node[midway, above left] {$\vec{r}$};
\end{tikzpicture}
\end{center}

Let $P$ be any point in space whose coordinates are:
\begin{itemize}[noitemsep]
    \item $(x, y, z)$ with respect to the old axes $OX, OY, OZ$.
    \item $(x', y', z')$ with respect to the new axes $O'X', O'Y', O'Z'$.
\end{itemize}

\begin{theorembox}[Transformation of Coordinates under Translation]
The relations connecting the old coordinates $(x,y,z)$ and the new coordinates $(x',y',z')$ are:
\begin{equation}
    \boxed{x = x' + h, \quad y = y' + k, \quad z = z' + l}
\end{equation}
Conversely, the new coordinates in terms of the old coordinates are:
\begin{equation}
    \boxed{x' = x - h, \quad y' = y - k, \quad z' = z - l}
\end{equation}
\end{theorembox}

\begin{proof}
Let the position vectors of $O'$ and $P$ with respect to the original origin $O$ be $\vec{r}_0$ and $\vec{r}$ respectively:
\begin{align}
    \vec{r}_0 &= h\,\hat{i} + k\,\hat{j} + l\,\hat{k}, \\
    \vec{r} &= x\,\hat{i} + y\,\hat{j} + z\,\hat{k}.
\end{align}
The position vector of $P$ relative to the new origin $O'$ is $\vec{r}' = x'\,\hat{i} + y'\,\hat{j} + z'\,\hat{k}$.
By vector addition in triangle $\triangle OO'P$:
\begin{equation}
    \vec{OP} = \vec{OO'} + \vec{O'P} \implies \vec{r} = \vec{r}_0 + \vec{r}'.
\end{equation}
Equating components along the three mutually orthogonal basis vectors $\hat{i}, \hat{j}, \hat{k}$:
\begin{equation}
    x = x' + h, \quad y = y' + k, \quad z = z' + l.
\end{equation}
\end{proof}

\subsection{Invariance of Spatial Distances}
A fundamental geometric invariant under pure translation is the Euclidean distance between any pair of points:
\begin{align}
    d^2 &= (x_2 - x_1)^2 + (y_2 - y_1)^2 + (z_2 - z_1)^2 \nonumber\\
    &= \bigl[(x'_2 + h) - (x'_1 + h)\bigr]^2 + \bigl[(y'_2 + k) - (y'_1 + k)\bigr]^2 + \bigl[(z'_2 + l) - (z'_1 + l)\bigr]^2 \nonumber\\
    &= (x'_2 - x'_1)^2 + (y'_2 - y'_1)^2 + (z'_2 - z'_1)^2.
\end{align}
Thus, distance is strictly invariant under translation of axes.

% =========================================================================
% PAGE 17
% =========================================================================
\section{Page 17: Projection of a Directed Line Segment on Coordinate Axes}
\label{sec:page17}

\begin{conceptbox}[Manuscript Overview: Page 17]
Page 17 analyzes the projection of a directed line segment $OP$ on the three rectangular coordinate axes. It establishes that the coordinates $(x,y,z)$ of any point $P$ are precisely the orthogonal projections of the position vector $\vec{OP}$ on the $x$-, $y$-, and $z$-axes.
\end{conceptbox}

\subsection{Geometric Definition of Orthogonal Projection}
The \textbf{orthogonal projection} of a directed line segment $AB$ onto a directed line $L$ is the directed segment $A'B'$ on $L$, where $A'$ and $B'$ are the feet of the perpendiculars dropped from $A$ and $B$ respectively onto $L$.
If $\theta$ is the angle between the directed segment $\vec{AB}$ and the directed line $L$, then the algebraic magnitude of the projection is:
\begin{equation}
    \text{Proj}_L(\vec{AB}) = |\vec{AB}|\cos\theta.
\end{equation}

\subsection{Projections of the Radius Vector $OP$}
Let $P(x,y,z)$ be any point in space, and let $OP = r = \sqrt{x^2+y^2+z^2}$.
Let the directed line segment $OP$ make angles $\alpha, \beta, \gamma$ with the positive directions of the coordinate axes $OX, OY, OZ$ respectively. These angles are called the \textbf{direction angles} of the line $OP$.

\begin{center}
\begin{tikzpicture}[scale=1.1, >=Stealth, 3d view={120}{25}]
    % Axes
    \draw[->, thick] (0,0,0) -- (4.2,0,0) node[anchor=north east]{$x$};
    \draw[->, thick] (0,0,0) -- (0,4.8,0) node[anchor=north west]{$y$};
    \draw[->, thick] (0,0,0) -- (0,0,3.8) node[anchor=south]{$z$};

    \coordinate (O) at (0,0,0);
    \coordinate (P) at (2.8, 3.6, 2.4);
    \coordinate (A) at (2.8,0,0);
    \coordinate (B) at (0,3.6,0);
    \coordinate (C) at (0,0,2.4);

    % Segment OP
    \draw[very thick, primary] (O) -- (P) node[midway, above left] {$r = OP$};
    \fill[accent] (P) circle (2pt) node[above right] {$P(x,y,z)$};

    % Projection right triangles
    \draw[dashed, secondary, thick] (P) -- (A) node[midway, below right] {$PA \perp OX$};
    \draw[dashed, secondary, thick] (P) -- (B);
    \draw[dashed, secondary, thick] (P) -- (C);

    \node[below] at (A) {$A(x,0,0)$};
    \node[left] at (B) {$B(0,y,0)$};
    \node[above left] at (C) {$C(0,0,z)$};
\end{tikzpicture}
\end{center}

In the right-angled triangle $\triangle OAP$, the angle at $A$ is $90^\circ$ because $PA$ lies in the plane perpendicular to the $x$-axis:
\begin{equation}
    \cos\alpha = \frac{OA}{OP} = \frac{x}{r} \implies \boxed{x = r\cos\alpha}
\end{equation}
Similarly, from the right-angled triangles $\triangle OBP$ and $\triangle OCP$:
\begin{align}
    \cos\beta &= \frac{OB}{OP} = \frac{y}{r} \implies \boxed{y = r\cos\beta} \\
    \cos\gamma &= \frac{OC}{OP} = \frac{z}{r} \implies \boxed{z = r\cos\gamma}
\end{align}

\begin{theorembox}[Coordinates as Projections]
The coordinates $(x,y,z)$ of any point $P$ at distance $r$ from the origin are the orthogonal projections of the radius vector $OP$ on the three coordinate axes:
\begin{equation}
    x = r\cos\alpha, \quad y = r\cos\beta, \quad z = r\cos\gamma
\end{equation}
where $\alpha, \beta, \gamma$ are the direction angles of $OP$.
\end{theorembox}

% =========================================================================
% PAGE 18
% =========================================================================
\section{Page 18: Projection of a Broken Line --- The Contour Theorem}
\label{sec:page18}

\begin{conceptbox}[Manuscript Overview: Page 18]
Page 18 derives the fundamental projection theorem for broken polygonal lines (the \textbf{Contour Theorem}). It proves that the projection of the straight segment joining the initial point to the terminal point equals the algebraic sum of the projections of the individual segments composing the path.
\end{conceptbox}

\subsection{The Fundamental Projection Theorem}
\begin{theorembox}[Projection of a Broken Line (Contour Theorem)]
Let $P_0, P_1, P_2, \dots, P_n$ be any sequence of points in space forming a broken line (polygonal path). Then the orthogonal projection of the direct segment $P_0 P_n$ on any directed line $L$ is equal to the algebraic sum of the projections of the individual directed segments $P_0 P_1, P_1 P_2, \dots, P_{n-1} P_n$ on $L$:
\begin{equation}
    \boxed{\text{Proj}_L(P_0 P_n) = \sum_{i=1}^n \text{Proj}_L(P_{i-1} P_i)}
\end{equation}
\end{theorembox}

\begin{proof}
Let $L$ be a directed line, and let $\hat{u}$ be a unit vector along $L$.
By vector addition, the displacement vector from $P_0$ to $P_n$ is:
\begin{equation}
    \vec{P_0 P_n} = \vec{P_0 P_1} + \vec{P_1 P_2} + \vec{P_2 P_3} + \dots + \vec{P_{n-1} P_n}.
\end{equation}
The orthogonal projection of any directed segment $\vec{AB}$ on $L$ is given by the scalar dot product with the unit vector $\hat{u}$:
\begin{equation}
    \text{Proj}_L(\vec{AB}) = \vec{AB} \cdot \hat{u}.
\end{equation}
Taking the dot product of both sides with $\hat{u}$:
\begin{align}
    \vec{P_0 P_n} \cdot \hat{u} &= \left( \sum_{i=1}^n \vec{P_{i-1} P_i} \right) \cdot \hat{u} \nonumber\\
    &= \sum_{i=1}^n (\vec{P_{i-1} P_i} \cdot \hat{u}).
\end{align}
Since $\vec{P_0 P_n} \cdot \hat{u} = \text{Proj}_L(P_0 P_n)$ and $\vec{P_{i-1} P_i} \cdot \hat{u} = \text{Proj}_L(P_{i-1} P_i)$, the theorem is proved.
\end{proof}

\subsection{Application: Projections of Arbitrary Spatial Displacements}
Consider two points $P(x_1, y_1, z_1)$ and $Q(x_2, y_2, z_2)$. We can go from $P$ to $Q$ along a broken path of three mutually orthogonal segments parallel to the coordinate axes:
\begin{enumerate}[noitemsep]
    \item From $P(x_1, y_1, z_1)$ to $A(x_2, y_1, z_1)$ parallel to the $x$-axis, length $\Delta x = x_2 - x_1$.
    \item From $A(x_2, y_1, z_1)$ to $B(x_2, y_2, z_1)$ parallel to the $y$-axis, length $\Delta y = y_2 - y_1$.
    \item From $B(x_2, y_2, z_1)$ to $Q(x_2, y_2, z_2)$ parallel to the $z$-axis, length $\Delta z = z_2 - z_1$.
\end{enumerate}
By the Contour Theorem, for any directed line $L$:
\begin{equation}
    \text{Proj}_L(PQ) = \text{Proj}_L(PA) + \text{Proj}_L(AB) + \text{Proj}_L(BQ).
\end{equation}
This theorem forms the operational foundation for calculating direction cosines, distances, and angles between lines in space.

% =========================================================================
% PAGE 19
% =========================================================================
\section{Page 19: Direction Cosines and the Fundamental Identity $l^2+m^2+n^2=1$}
\label{sec:page19}

\begin{conceptbox}[Manuscript Overview: Page 19]
Page 19 formalizes the definitions of the direction cosines $l = \cos\alpha$, $m = \cos\beta$, and $n = \cos\gamma$. It proves the fundamental quadratic identity $l^2 + m^2 + n^2 = 1$ using the rectangular parallelepiped geometry and establishes key trigonometric corollaries.
\end{conceptbox}

\subsection{Definition of Direction Cosines}
\begin{conceptbox}[Direction Cosines (DCs)]
If a directed straight line makes angles $\alpha, \beta, \gamma$ with the positive directions of the coordinate axes $OX, OY, OZ$ respectively (where $0 \le \alpha, \beta, \gamma \le \pi$), then the cosines of these angles are called the \textbf{direction cosines} of the line, denoted conventionally by:
\begin{equation}
    l = \cos\alpha, \quad m = \cos\beta, \quad n = \cos\gamma.
\end{equation}
\end{conceptbox}

\subsection{The Fundamental Quadratic Identity}
\begin{theorembox}[Fundamental Identity for Direction Cosines]
For any straight line in three-dimensional space with direction cosines $(l, m, n)$:
\begin{equation}
    \boxed{l^2 + m^2 + n^2 = 1 \iff \cos^2\alpha + \cos^2\beta + \cos^2\gamma = 1}
\end{equation}
\end{theorembox}

\begin{proof}
Let a line through the origin $O$ have direction cosines $l, m, n$.
Let $P(x, y, z)$ be any point on this line at distance $r$ from $O$, so that $OP = r > 0$.
From the projection relations established on Page 17:
\begin{equation}
    x = r\cos\alpha = lr, \quad y = r\cos\beta = mr, \quad z = r\cos\gamma = nr.
\end{equation}
From the 3D distance formula, the squared distance $OP^2$ is:
\begin{equation}
    r^2 = x^2 + y^2 + z^2.
\end{equation}
Substituting the expressions for $x, y, z$:
\begin{align}
    r^2 &= (lr)^2 + (mr)^2 + (nr)^2 \nonumber\\
    r^2 &= r^2(l^2 + m^2 + n^2).
\end{align}
Since $r > 0$, we divide both sides by $r^2$:
\begin{equation}
    \boxed{l^2 + m^2 + n^2 = 1}
\end{equation}
\end{proof}

\begin{notebox}[Trigonometric Corollaries]
\begin{enumerate}
    \item \textbf{Sum of Squared Sines}:
    Using $\sin^2\theta = 1 - \cos^2\theta$ for each angle:
    \begin{align}
        \sin^2\alpha + \sin^2\beta + \sin^2\gamma &= (1 - \cos^2\alpha) + (1 - \cos^2\beta) + (1 - \cos^2\gamma) \nonumber\\
        &= 3 - (\cos^2\alpha + \cos^2\beta + \cos^2\gamma) = 3 - 1 = \boxed{2}.
    \end{align}
    \item \textbf{Sum of Double Angles}:
    Using $\cos 2\theta = 2\cos^2\theta - 1$:
    \begin{align}
        \cos 2\alpha + \cos 2\beta + \cos 2\gamma &= (2\cos^2\alpha - 1) + (2\cos^2\beta - 1) + (2\cos^2\gamma - 1) \nonumber\\
        &= 2(\cos^2\alpha + \cos^2\beta + \cos^2\gamma) - 3 = 2(1) - 3 = \boxed{-1}.
    \end{align}
    \item \textbf{Range of DCs}: Since $l^2 \le l^2+m^2+n^2=1$, every direction cosine satisfies:
    \begin{equation}
        -1 \le l \le 1, \quad -1 \le m \le 1, \quad -1 \le n \le 1.
    \end{equation}
\end{enumerate}
\end{notebox}

% =========================================================================
% PAGE 20
% =========================================================================
\section{Page 20: Projection of Segment $PQ$ on a Line with Direction Cosines $(l,m,n)$}
\label{sec:page20}

\begin{conceptbox}[Manuscript Overview: Page 20]
Page 20 derives the general formula for the orthogonal projection of a segment joining $P(x_1, y_1, z_1)$ and $Q(x_2, y_2, z_2)$ onto an arbitrary straight line in space having direction cosines $(l, m, n)$.
\end{conceptbox}

\subsection{Step-by-Step Derivation of the Projection Formula}
\begin{theorembox}[Projection of a Segment on an Arbitrary Line]
The orthogonal projection of the directed line segment joining $P(x_1, y_1, z_1)$ and $Q(x_2, y_2, z_2)$ onto a line $L$ whose direction cosines are $(l, m, n)$ is given by:
\begin{equation}
    \boxed{\text{Proj}_L(PQ) = l(x_2 - x_1) + m(y_2 - y_1) + n(z_2 - z_1)}
\end{equation}
\end{theorembox}

\begin{proof}
Let $L$ be a line passing through the origin (without loss of generality, since projection depends only on direction) having direction cosines $(l, m, n) = (\cos\alpha, \cos\beta, \cos\gamma)$.
The unit vector along line $L$ is:
\begin{equation}
    \hat{u} = l\,\hat{i} + m\,\hat{j} + n\,\hat{k}.
\end{equation}
The displacement vector from $P$ to $Q$ is:
\begin{equation}
    \vec{PQ} = (x_2 - x_1)\,\hat{i} + (y_2 - y_1)\,\hat{j} + (z_2 - z_1)\,\hat{k}.
\end{equation}
By the geometric definition of projection via scalar product:
\begin{align}
    \text{Proj}_L(PQ) &= \vec{PQ} \cdot \hat{u} \nonumber\\
    &= \bigl[(x_2 - x_1)\,\hat{i} + (y_2 - y_1)\,\hat{j} + (z_2 - z_1)\,\hat{k}\bigr] \cdot (l\,\hat{i} + m\,\hat{j} + n\,\hat{k}) \nonumber\\
    &= l(x_2 - x_1) + m(y_2 - y_1) + n(z_2 - z_1).
\end{align}
Alternatively, using the Contour Theorem of Page 18:
Let the coordinate displacements be $PA = x_2 - x_1$ (parallel to $OX$), $AB = y_2 - y_1$ (parallel to $OY$), and $BQ = z_2 - z_1$ (parallel to $OZ$).
The angle between $PA$ and line $L$ is $\alpha$, so $\text{Proj}_L(PA) = (x_2 - x_1)\cos\alpha = l(x_2 - x_1)$.
The angle between $AB$ and line $L$ is $\beta$, so $\text{Proj}_L(AB) = (y_2 - y_1)\cos\beta = m(y_2 - y_1)$.
The angle between $BQ$ and line $L$ is $\gamma$, so $\text{Proj}_L(BQ) = (z_2 - z_1)\cos\gamma = n(z_2 - z_1)$.
Summing the three projections yields:
\begin{equation}
    \text{Proj}_L(PQ) = l(x_2 - x_1) + m(y_2 - y_1) + n(z_2 - z_1).
\end{equation}
\end{proof}

\subsection{Derivation of the Distance Formula via Self-Projection}
If line $L$ is chosen along the directed segment $PQ$ itself, then the projection of $PQ$ onto $L$ is its full length $d = PQ$.
The direction cosines of the segment $PQ$ are:
\begin{equation}
    l = \frac{x_2 - x_1}{d}, \quad m = \frac{y_2 - y_1}{d}, \quad n = \frac{z_2 - z_1}{d}.
\end{equation}
Applying the projection formula:
\begin{align}
    d &= l(x_2 - x_1) + m(y_2 - y_1) + n(z_2 - z_1) \nonumber\\
    &= \frac{x_2 - x_1}{d}(x_2 - x_1) + \frac{y_2 - y_1}{d}(y_2 - y_1) + \frac{z_2 - z_1}{d}(z_2 - z_1) \nonumber\\
    d^2 &= (x_2 - x_1)^2 + (y_2 - y_1)^2 + (z_2 - z_1)^2 \nonumber\\
    \implies d &= \sqrt{(x_2 - x_1)^2 + (y_2 - y_1)^2 + (z_2 - z_1)^2}.
\end{align}

% =========================================================================
% PAGE 21
% =========================================================================
\section{Page 21: Direction Ratios and Normalization to Direction Cosines}
\label{sec:page21}

\begin{conceptbox}[Manuscript Overview: Page 21]
Page 21 introduces \textbf{Direction Ratios} (DRs), which avoid the quadratic constraint $l^2+m^2+n^2=1$. It establishes the normalization algorithm to recover direction cosines from ratios and derives the analytical condition for the perpendicularity of two lines.
\end{conceptbox}

\subsection{Definition of Direction Ratios}
\begin{conceptbox}[Direction Ratios (DRs)]
Any three real numbers $a, b, c$ that are proportional to the direction cosines $l, m, n$ of a line are called the \textbf{direction ratios} (or direction numbers) of that line:
\begin{equation}
    \frac{l}{a} = \frac{m}{b} = \frac{n}{c} = k \quad (k \neq 0).
\end{equation}
\end{conceptbox}

\subsection{Normalization Formula}
\begin{theorembox}[Conversion from Direction Ratios to Direction Cosines]
If $a, b, c$ are the direction ratios of a line, then its actual direction cosines $l, m, n$ are given by:
\begin{equation}
    \boxed{l = \pm \frac{a}{\sqrt{a^2 + b^2 + c^2}}, \quad m = \pm \frac{b}{\sqrt{a^2 + b^2 + c^2}}, \quad n = \pm \frac{c}{\sqrt{a^2 + b^2 + c^2}}}
\end{equation}
where either all upper signs or all lower signs are taken simultaneously.
\end{theorembox}

\begin{proof}
Let $\frac{l}{a} = \frac{m}{b} = \frac{n}{c} = k$. Then:
\begin{equation}
    l = ka, \quad m = kb, \quad n = kc.
\end{equation}
Squaring and adding these equations:
\begin{equation}
    l^2 + m^2 + n^2 = k^2(a^2 + b^2 + c^2).
\end{equation}
Since $l^2 + m^2 + n^2 = 1$:
\begin{equation}
    1 = k^2(a^2 + b^2 + c^2) \implies k = \pm \frac{1}{\sqrt{a^2 + b^2 + c^2}}.
\end{equation}
Substituting $k$ back into $l = ka, m = kb, n = kc$ yields the normalization equations.
\end{proof}

\subsection{Direction Ratios of the Line Joining Two Points}
For a line passing through $P(x_1, y_1, z_1)$ and $Q(x_2, y_2, z_2)$, its direction cosines are:
\begin{equation}
    l = \frac{x_2 - x_1}{d}, \quad m = \frac{y_2 - y_1}{d}, \quad n = \frac{z_2 - z_1}{d}, \quad \text{where } d = PQ.
\end{equation}
Multiplying each by $d$ demonstrates that:
\begin{equation}
    \boxed{(a, b, c) = (x_2 - x_1, \; y_2 - y_1, \; z_2 - z_1)}
\end{equation}
are a valid set of direction ratios for the line joining $P$ and $Q$.

\subsection{Condition for Perpendicularity of Two Lines}
Two lines $L_1$ and $L_2$ with direction cosines $(l_1, m_1, n_1)$ and $(l_2, m_2, n_2)$ are perpendicular ($\theta = 90^\circ \iff \cos\theta = 0$) if and only if:
\begin{equation}
    \boxed{l_1 l_2 + m_1 m_2 + n_1 n_2 = 0}
\end{equation}
In terms of direction ratios $(a_1, b_1, c_1)$ and $(a_2, b_2, c_2)$:
\begin{equation}
    \boxed{a_1 a_2 + b_1 b_2 + c_1 c_2 = 0}
\end{equation}

% =========================================================================
% PAGE 22
% =========================================================================
\section{Page 22: Angle Between Two Directed Lines in Space}
\label{sec:page22}

\begin{conceptbox}[Manuscript Overview: Page 22]
Page 22 derives the cosine formula for the angle $\theta$ between two straight lines in three dimensions, given their direction cosines $(l_1, m_1, n_1)$ and $(l_2, m_2, n_2)$ or their direction ratios.
\end{conceptbox}

\subsection{Derivation of the Cosine Formula}
\begin{theorembox}[Angle Between Two Lines]
Let $L_1$ and $L_2$ be two straight lines in space having direction cosines $(l_1, m_1, n_1)$ and $(l_2, m_2, n_2)$. The angle $\theta$ between them is given by:
\begin{equation}
    \boxed{\cos\theta = l_1 l_2 + m_1 m_2 + n_1 n_2}
\end{equation}
In terms of direction ratios $(a_1, b_1, c_1)$ and $(a_2, b_2, c_2)$:
\begin{equation}
    \boxed{\cos\theta = \frac{a_1 a_2 + b_1 b_2 + c_1 c_2}{\sqrt{a_1^2 + b_1^2 + c_1^2}\,\sqrt{a_2^2 + b_2^2 + c_2^2}}}
\end{equation}
\end{theorembox}

\begin{proof}
Draw two lines through the origin $O$ parallel to $L_1$ and $L_2$.
On these lines, choose points $P$ and $Q$ such that $OP = 1$ and $OQ = 1$.
The coordinates of $P$ and $Q$ are:
\begin{equation}
    P(l_1, m_1, n_1), \quad Q(l_2, m_2, n_2).
\end{equation}
In triangle $\triangle OPQ$, let $\angle POQ = \theta$.
By the law of cosines in $\triangle OPQ$:
\begin{equation}
    PQ^2 = OP^2 + OQ^2 - 2(OP)(OQ)\cos\theta = 1^2 + 1^2 - 2(1)(1)\cos\theta = 2 - 2\cos\theta.
\end{equation}
Using the 3D distance formula to evaluate $PQ^2$:
\begin{align}
    PQ^2 &= (l_2 - l_1)^2 + (m_2 - m_1)^2 + (n_2 - n_1)^2 \nonumber\\
    &= (l_2^2 + m_2^2 + n_2^2) + (l_1^2 + m_1^2 + n_1^2) - 2(l_1 l_2 + m_1 m_2 + n_1 n_2).
\end{align}
Since $l_1^2+m_1^2+n_1^2 = 1$ and $l_2^2+m_2^2+n_2^2 = 1$:
\begin{equation}
    PQ^2 = 1 + 1 - 2(l_1 l_2 + m_1 m_2 + n_1 n_2) = 2 - 2(l_1 l_2 + m_1 m_2 + n_1 n_2).
\end{equation}
Equating the two expressions for $PQ^2$:
\begin{equation}
    2 - 2\cos\theta = 2 - 2(l_1 l_2 + m_1 m_2 + n_1 n_2) \implies \boxed{\cos\theta = l_1 l_2 + m_1 m_2 + n_1 n_2}.
\end{equation}
Substituting $l_i = \frac{a_i}{\sqrt{\sum a_i^2}}$ gives the direction ratio formulation.
\end{proof}

\begin{center}
\begin{tikzpicture}[scale=1.2, >=Stealth]
    \coordinate (O) at (0,0);
    \coordinate (P) at (25:3.2);
    \coordinate (Q) at (75:3.2);

    \draw[->, very thick, primary] (O) -- (P) node[right] {$P(l_1, m_1, n_1)$};
    \draw[->, very thick, secondary] (O) -- (Q) node[above] {$Q(l_2, m_2, n_2)$};
    \draw[thick, accent, dashed] (P) -- (Q) node[midway, above right] {$PQ$};

    \draw[->, thick] (1,0) arc (0:25:1) node[midway, right=1pt] {$\theta_1$};
    \draw[->, thick, primary] (0.8, {0.8*tan(25)}) arc (25:75:0.8) node[midway, above right] {$\theta$};

    \node[below left] at (O) {$O$};
    \node[midway, below] at ($(O)!0.5!(P)$) {$OP = 1$};
    \node[midway, left] at ($(O)!0.5!(Q)$) {$OQ = 1$};
\end{tikzpicture}
\end{center}

% =========================================================================
% PAGE 23
% =========================================================================
\section{Page 23: Sine of the Angle, Lagrange's Identity, and Parallelism}
\label{sec:page23}

\begin{conceptbox}[Manuscript Overview: Page 23]
Page 23 derives the formula for $\sin\theta$ between two lines using Lagrange's algebraic identity and establishes the condition for two lines to be parallel in space.
\end{conceptbox}

\subsection{Lagrange's Identity and the Sine Formula}
\begin{theorembox}[Sine of the Angle Between Two Lines]
The sine of the angle $\theta$ between two lines with direction cosines $(l_1, m_1, n_1)$ and $(l_2, m_2, n_2)$ is given by:
\begin{equation}
    \boxed{\sin\theta = \sqrt{(l_1 m_2 - l_2 m_1)^2 + (m_1 n_2 - m_2 n_1)^2 + (n_1 l_2 - n_2 l_1)^2}}
\end{equation}
\end{theorembox}

\begin{proof}
Using $\sin^2\theta = 1 - \cos^2\theta$:
Since $l_1^2+m_1^2+n_1^2 = 1$ and $l_2^2+m_2^2+n_2^2 = 1$, we can write $1$ as their product:
\begin{equation}
    \sin^2\theta = (l_1^2 + m_1^2 + n_1^2)(l_2^2 + m_2^2 + n_2^2) - (l_1 l_2 + m_1 m_2 + n_1 n_2)^2.
\end{equation}
By \textbf{Lagrange's Algebraic Identity}:
For any six real numbers $(A, B, C)$ and $(X, Y, Z)$:
\begin{equation}
    (A^2 + B^2 + C^2)(X^2 + Y^2 + Z^2) - (AX + BY + CZ)^2 = (AY - BX)^2 + (BZ - CY)^2 + (CX - AZ)^2.
\end{equation}
Setting $(A, B, C) = (l_1, m_1, n_1)$ and $(X, Y, Z) = (l_2, m_2, n_2)$:
\begin{equation}
    \sin^2\theta = (l_1 m_2 - l_2 m_1)^2 + (m_1 n_2 - m_2 n_1)^2 + (n_1 l_2 - n_2 l_1)^2.
\end{equation}
Taking the principal square root (since $0 \le \theta \le \pi \implies \sin\theta \ge 0$):
\begin{equation}
    \sin\theta = \sqrt{(l_1 m_2 - l_2 m_1)^2 + (m_1 n_2 - m_2 n_1)^2 + (n_1 l_2 - n_2 l_1)^2}.
\end{equation}
\end{proof}

\subsection{Condition for Parallelism of Two Lines}
Two lines $L_1$ and $L_2$ are parallel if and only if the angle between them is $\theta = 0$ or $\theta = \pi$, which implies $\sin\theta = 0$.
Since $\sin^2\theta$ is the sum of three non-negative squared terms:
\begin{equation}
    (l_1 m_2 - l_2 m_1)^2 + (m_1 n_2 - m_2 n_1)^2 + (n_1 l_2 - n_2 l_1)^2 = 0.
\end{equation}
A sum of squares of real numbers vanishes if and only if each individual term vanishes simultaneously:
\begin{equation}
    l_1 m_2 - l_2 m_1 = 0, \quad m_1 n_2 - m_2 n_1 = 0, \quad n_1 l_2 - n_2 l_1 = 0.
\end{equation}
This yields the celebrated proportionality condition:
\begin{equation}
    \boxed{\frac{l_1}{l_2} = \frac{m_1}{m_2} = \frac{n_1}{n_2}}
\end{equation}
In terms of direction ratios $(a_1, b_1, c_1)$ and $(a_2, b_2, c_2)$:
\begin{equation}
    \boxed{\frac{a_1}{a_2} = \frac{b_1}{b_2} = \frac{c_1}{c_2}}
\end{equation}

% =========================================================================
% PAGE 24
% =========================================================================
\section{Page 24: Solved Problem --- The Four Diagonals of a Cube}
\label{sec:page24}

\begin{conceptbox}[Manuscript Overview: Page 24]
Page 24 presents a classic examination problem: determining the direction cosines of the four body diagonals of a cube and computing the angles between any pair of diagonals, as well as between a diagonal and an edge.
\end{conceptbox}

\begin{examplebox}[Examination Problem: The Diagonals of a Cube]
Prove that the angle between any two diagonals of a cube is $\cos^{-1}\left(\frac{1}{3}\right) \approx 70^\circ 32'$.
\end{examplebox}

\begin{proof}[Complete Step-by-Step Solution]
Let $OABC-DEFG$ be a cube of edge length $a$.
Set the origin at vertex $O(0,0,0)$, and let the three coterminous edges $OA, OC, OD$ lie along the positive directions of the $x$-, $y$-, and $z$-axes respectively.

\begin{center}
\begin{tikzpicture}[scale=1.2, >=Stealth, 3d view={125}{25}]
    \def\a{3.2}
    % Axes
    \draw[->, thick] (0,0,0) -- (4.2,0,0) node[anchor=north east]{$x$};
    \draw[->, thick] (0,0,0) -- (0,4.2,0) node[anchor=north west]{$y$};
    \draw[->, thick] (0,0,0) -- (0,0,4.2) node[anchor=south]{$z$};

    % Vertices
    \coordinate (O) at (0,0,0);
    \coordinate (A) at (\a,0,0);
    \coordinate (C) at (0,\a,0);
    \coordinate (D) at (0,0,\a);
    \coordinate (B) at (\a,\a,0);
    \coordinate (E) at (\a,0,\a);
    \coordinate (G) at (0,\a,\a);
    \coordinate (F) at (\a,\a,\a);

    % Edges
    \draw[thick] (A) -- (B) -- (C);
    \draw[thick] (A) -- (E) -- (F) -- (B);
    \draw[thick] (C) -- (G) -- (F);
    \draw[thick] (D) -- (E);
    \draw[thick] (D) -- (G);
    \draw[dashed, thick] (O) -- (A);
    \draw[dashed, thick] (O) -- (C);
    \draw[dashed, thick] (O) -- (D);

    % Diagonals
    \draw[very thick, primary] (O) -- (F) node[midway, above left] {$OF$};
    \draw[very thick, accent, dashed] (A) -- (G) node[midway, below right] {$AG$};

    % Nodes
    \node[below left] at (O) {$O(0,0,0)$};
    \node[below] at (A) {$A(a,0,0)$};
    \node[right] at (C) {$C(0,a,0)$};
    \node[left] at (D) {$D(0,0,a)$};
    \node[above right] at (F) {$F(a,a,a)$};
    \node[above left] at (G) {$G(0,a,a)$};
\end{tikzpicture}
\end{center}

The coordinates of the $8$ vertices of the cube are:
\begin{align}
    &O(0,0,0), \quad A(a,0,0), \quad C(0,a,0), \quad D(0,0,a), \nonumber\\
    &B(a,a,0), \quad E(a,0,a), \quad G(0,a,a), \quad F(a,a,a).
\end{align}
The cube possesses $4$ body diagonals connecting opposite pairs of vertices:
\begin{enumerate}
    \item Diagonal $OF$: joins $O(0,0,0)$ to $F(a,a,a)$.
    \begin{align}
        \text{DRs} &= (a - 0, \; a - 0, \; a - 0) = (a, a, a) \sim (1, 1, 1). \nonumber\\
        \text{Length } OF &= \sqrt{a^2 + a^2 + a^2} = a\sqrt{3}. \nonumber\\
        \text{DCs of } OF &= \left(\frac{1}{\sqrt{3}}, \; \frac{1}{\sqrt{3}}, \; \frac{1}{\sqrt{3}}\right).
    \end{align}
    \item Diagonal $AG$: joins $A(a,0,0)$ to $G(0,a,a)$.
    \begin{align}
        \text{DRs} &= (0 - a, \; a - 0, \; a - 0) = (-a, a, a) \sim (-1, 1, 1). \nonumber\\
        \text{DCs of } AG &= \left(-\frac{1}{\sqrt{3}}, \; \frac{1}{\sqrt{3}}, \; \frac{1}{\sqrt{3}}\right).
    \end{align}
    \item Diagonal $BE$: joins $B(a,a,0)$ to $E(0,0,a)$ (or $(a,0,a)$ to $(0,a,0)$):
    \begin{align}
        \text{DRs} &\sim (1, -1, 1). \nonumber\\
        \text{DCs of } BE &= \left(\frac{1}{\sqrt{3}}, \; -\frac{1}{\sqrt{3}}, \; \frac{1}{\sqrt{3}}\right).
    \end{align}
    \item Diagonal $CD$: joins $C(0,a,0)$ to $D(a,a,0)$'s opposite vertex:
    \begin{align}
        \text{DRs} &\sim (1, 1, -1). \nonumber\\
        \text{DCs of } CD &= \left(\frac{1}{\sqrt{3}}, \; \frac{1}{\sqrt{3}}, \; -\frac{1}{\sqrt{3}}\right).
    \end{align}
\end{enumerate}

Now, let $\theta$ be the angle between any two diagonals, say $OF$ and $AG$:
\begin{align}
    \cos\theta &= l_1 l_2 + m_1 m_2 + n_1 n_2 \nonumber\\
    &= \left(\frac{1}{\sqrt{3}}\right)\left(-\frac{1}{\sqrt{3}}\right) + \left(\frac{1}{\sqrt{3}}\right)\left(\frac{1}{\sqrt{3}}\right) + \left(\frac{1}{\sqrt{3}}\right)\left(\frac{1}{\sqrt{3}}\right) \nonumber\\
    &= -\frac{1}{3} + \frac{1}{3} + \frac{1}{3} = \frac{1}{3}.
\end{align}
For any other pair of diagonals, the calculation yields $\pm \frac{1}{3}$. Taking the acute angle between the diagonals:
\begin{equation}
    \boxed{\theta = \cos^{-1}\left(\frac{1}{3}\right) \approx 70^\circ 31' 44''}
\end{equation}
\end{proof}

\begin{notebox}[Angle Between a Body Diagonal and an Edge]
Consider the diagonal $OF$ with DCs $\left(\frac{1}{\sqrt{3}}, \frac{1}{\sqrt{3}}, \frac{1}{\sqrt{3}}\right)$ and the edge $OA$ lying along the $x$-axis with DCs $(1, 0, 0)$.
The angle $\phi$ between them is:
\begin{equation}
    \cos\phi = \left(\frac{1}{\sqrt{3}}\right)(1) + \left(\frac{1}{\sqrt{3}}\right)(0) + \left(\frac{1}{\sqrt{3}}\right)(0) = \frac{1}{\sqrt{3}} \implies \phi = \cos^{-1}\left(\frac{1}{\sqrt{3}}\right) \approx 54^\circ 44'.
\end{equation}
Similarly, the angle $\psi$ between a body diagonal and a face diagonal (e.g., $OB$ with DCs $\left(\frac{1}{\sqrt{2}}, \frac{1}{\sqrt{2}}, 0\right)$) is:
\begin{equation}
    \cos\psi = \left(\frac{1}{\sqrt{3}}\right)\left(\frac{1}{\sqrt{2}}\right) + \left(\frac{1}{\sqrt{3}}\right)\left(\frac{1}{\sqrt{2}}\right) + 0 = \frac{2}{\sqrt{6}} = \sqrt{\frac{2}{3}} \implies \psi = \cos^{-1}\left(\sqrt{\frac{2}{3}}\right) \approx 35^\circ 16'.
\end{equation}
\end{notebox}
